\ifdefined\gkver\else\def\gkver{student}\fi
\documentclass[\gkver]{gaokaozhenti}
\begin{document}

\chapter{2026年河北}

\begin{enumerate}
\item
%% number: 1
%% paperName: 2026年河北高考第 1 题 4 分
%% typeId: 1
%% type: 单选题
%% chapter: 直线运动
%% point: 路程与位移
%% method: 无
%% score: 4
%% degree: 950
%% duplicateId: 0
%% body:
某同学健身时，哑铃被举高 $0.5\Um$ 后再回到原位置，如图所示。如此重复 10 次，哑铃的位移大小为\xzanswer{A}
\onepicture[w=2.4cm, align=right]{\includesvg{figs/01}}

\fourchoices[answer=A]
{$0$}
{$5\Um$}
{$10\Um$}
{$11\Um$}

%% answer: A
%% memo:
\memoanswer{
位移是描述物体位置变化的物理量，位移的大小等于初位置到末位置的有向线段长度，和运动过程的路径无关。

本题中，哑铃每次运动都从原位置出发，举高后回到原位置，重复 $10$ 次后，哑铃的末位置与初始位置重合，初末位置的距离为 $0$，因此位移大小为 $0$。

故选 A。
}

\item
%% number: 2
%% paperName: 2026年河北高考第 2 题 4 分
%% typeId: 1
%% type: 单选题
%% chapter: 电场
%% point: 电场线
%% method: 无
%% score: 4
%% degree: 850
%% duplicateId: 0
%% body:
河北雄安新区融合静电探测技术构建起“天空之网”，实现了对低空飞行器的全天候实时监测。当飞行器靠近时，固定在地面上的静电探测器探头感应出了正电荷。在感应电荷产生的静电场中，电场线\xzanswer{B}

\fourchoices[answer=B]
{是闭合曲线}
{从感应正电荷出发}
{在无电荷区域相交}
{终止于感应正电荷}

%% answer: B
%% memo:
\memoanswer{
本题考查静电场电场线的基本性质，结合性质逐一判断选项：

A．感应电荷产生的静电场的电场线为非闭合曲线，故 A 错误；

B．静电场电场线起始于正电荷或无穷远，探测器探头感应出正电荷，因此电场线从感应正电荷出发，故 B 正确；

C．电场中某点的场强方向唯一，因此无电荷区域电场线不可能相交，故 C 错误；

D．静电场电场线终止于负电荷或无穷远，正电荷是电场线的起点而非终点，故 D 错误。

故选 B。
}

\item
%% number: 3
%% paperName: 2026年河北高考第 3 题 4 分
%% typeId: 1
%% type: 单选题
%% chapter: 气体性质
%% point: 理想气体状态方程
%% method: 无
%% score: 4
%% degree: 750
%% duplicateId: 0
%% body:
一款自动监测环境气压的装置由粗细均匀、导热良好的 U 形管和自动监测系统组成。如图所示，竖直放置的 U 形管右端开口，左侧管内用水银封闭了一定质量、可视为理想气体的空气。监测系统通过监测右侧液面的变化。结合温度即可得到环境气压。若外界温度下降的同时环境气压也发生了变化，导致 U 形管右侧液面上升，忽略水银体积变化，则封闭气体的\xzanswer{C}
\onepicture[w=4.2cm, align=right]{\includesvg{figs/03}}

\fourchoices[answer=C]
{内能增加}
{压强增大}
{分子的数密度减少}
{分子热运动的平均动能不变}

%% answer: C
%% memo:
\memoanswer{
A．U 形管导热良好，外界温度下降，因此封闭气体的温度也会下降。理想气体的内能仅由温度决定，温度降低，内能减小，A 错误；

C．U 形管水银总体积不变，右侧液面上升，说明左侧封闭气体的体积 $V$ 增大，封闭气体分子总数不变，因此单位体积的分子数（分子数密度）减少，C 正确；

根据理想气体状态方程 $\frac{pV}{T}=C$，变形得 $p = C \frac{T}{V}$。

已知温度 $T$ 减小、体积 $V$ 增大，则封闭气体压强减小，B 错误；

D．温度降低，则分子热运动的平均动能减小，D 错误。

故选 C。
}

\item
%% number: 4
%% paperName: 2026年河北高考第 4 题 4 分
%% typeId: 1
%% type: 单选题
%% chapter: 功和能
%% point: 交通工具启动
%% method: 无
%% score: 4
%% degree: 750
%% duplicateId: 0
%% body:
我国科技爱好者复原了春秋战国时期带有刃车軎（wéi）的马车；并对其性能进行了测试。在 $t_{1}-t_{2}$ 时间内，若马车以恒定功率在水平路面上沿直线运动，速度从 $v_{1}$ 加速到 $v_{2}$，假定马车所受阻力不变，则马车运动的 $v-t$ 图像可能是\xzanswer{C}
\fourchoices[answer=C, ispicture=true, h=2.6cm]
{figs/04a.png}
{figs/04b.png}
{figs/04c.png}
{figs/04d.png}

%% answer: C
%% memo:
\memoanswer{
马车以恒定功率运动，阻力 $f$ 不变，由功率公式 $P = Fv$

结合牛顿第二定律 $F - f = ma$

整理得加速度 $a=\frac{P}{mv}-\frac{f}{m}$

随着速度 $v$ 增大，恒定功率 $P$ 不变，因此加速度 $a$ 逐渐减小。

$v-t$ 图像的斜率表示加速度，因此 $v-t$ 图的斜率应逐渐减小，图像越来越平缓。

故选 C。
}

\item
%% number: 5
%% paperName: 2026年河北高考第 5 题 4 分
%% typeId: 1
%% type: 单选题
%% chapter: 光学
%% point: 光的干涉
%% method: 无
%% score: 4
%% degree: 750
%% duplicateId: 0
%% body:
如图所示，一根粗细均匀的细金属丝置于两挡板之间的水平狭缝中央。其中心轴与狭缝平行，直径小于狭缝宽度。一束平行单色光垂直挡板入射，在屏幕上形成了稳定的双缝干涉图样，相邻两条亮条纹中心间距为 $\Delta x$。若金属丝经加热均匀膨胀，则\xzanswer{A}
\onepicture[w=5.0cm, align=right]{\includesvg{figs/05}}

\fourchoices[answer=A]
{$\Delta x$ 变小}
{$\Delta x$ 变大}
{双缝间的距离减小}
{双缝间的距离不变}

%% answer: A
%% memo:
\memoanswer{
CD．金属丝加热均匀膨胀，其直径 $D$ 变大。该装置利用金属丝两侧的空隙形成双缝干涉，设两挡板间狭缝总宽度为 $W$（不变），则双缝间距（即两透光狭缝中心间的距离）为 $d = \frac{W + D}{2}$

当 $D$ 变大时，双缝间距 $d$ 变大，故 CD 错误；

AB．根据双缝干涉条纹间距公式 $\Delta x = \frac{L}{d} \lambda$（其中 $L$ 为缝到屏的距离，$\lambda$ 为光波长），当 $d$ 变大时，$\Delta x$ 变小，故 A 正确，B 错误。

故选 A。
}

\item
%% number: 6
%% paperName: 2026年河北高考第 6 题 4 分
%% typeId: 1
%% type: 单选题
%% chapter: 电磁感应
%% point: 法拉第电磁感应定律
%% method: 无
%% score: 4
%% degree: 650
%% duplicateId: 0
%% body:
如图所示，在长直螺线管中间区域放置一个匝数为 $n$、面积为 $S$ 的同轴小线圈。已知长直螺线管通电后其内部磁场处处相同，磁感应强度的大小 $B$ 与螺线管中的电流 $I$ 成正比，即 $B = \alpha I$。在一段时间内，若电流 $I$ 随时间 $t$ 的变化关系满足 $I = \beta t$（$\beta$ 为常量），由理想电压表测出小线圈的感应电动势为 $E$，则 $\alpha$ 可以表示为\xzanswer{D}
\onepicture[w=7.5cm]{figs/06.png}

\fourchoices[answer=D]
{$\frac{\sqrt{2}E}{\beta S}$}
{$\frac{\sqrt{2}E}{n\beta S}$}
{$\frac{E}{\beta S}$}
{$\frac{E}{n\beta S}$}

%% answer: D
%% memo:
\memoanswer{
根据法拉第电磁感应定律有 $E = nS \cdot \frac{\Delta B}{\Delta t}$

根据题意有 $B = \alpha I$，$I = \beta t$

联立可得 $E = nS\alpha\beta$

解得 $\alpha = \frac{E}{n\beta S}$

故选 D。
}

\item
%% number: 7
%% paperName: 2026年河北高考第 7 题 4 分
%% typeId: 1
%% type: 单选题
%% chapter: 曲线运动
%% point: 人造地球卫星
%% method: 无
%% score: 4
%% degree: 350
%% duplicateId: 0
%% body:
为清除太空碎片对航天器的潜在威胁，某兴趣小组提出一种设想。如图所示，一质量为 $m$ 的太空碎片绕地球做半径为 $r$ 的匀速圆周运动，在 $Q$ 点受到一个与其速度方向垂直且背离地心向外的瞬时冲量 $I$ 作用、变轨到图中的椭圆轨道，最终进入大气层而烧毁，设地球质量为 $M$，引力常量为 $G$，则该太空碎片受到冲量作用后瞬间的动能为\xzanswer{B}
\onepicture[w=4.0cm, align=right]{\includesvg{figs/07}}

\fourchoices[answer=B]
{$\frac{GMm}{r} + \frac{I^2}{2m}$}
{$\frac{GMm}{2r} + \frac{I^2}{2m}$}
{$\frac{GMm}{r} + \frac{I^2}{2m} + I\sqrt{\frac{GM}{r}}$}
{$\frac{GMm}{2r} + \frac{I^2}{2m} + I\sqrt{\frac{GM}{r}}$}

%% answer: B
%% memo:
\memoanswer{
设太空碎片做匀速圆周运动的速度为 $v_0$，根据万有引力提供向心力有 $G\frac{Mm}{r^2} = m\frac{v_0^2}{r}$

解得 $v_0^2 = \frac{GM}{r}$

则碎片受到冲量前的初始动能为 $E_{k0} = \frac{1}{2}mv_0^2 = \frac{GMm}{2r}$

碎片在 $Q$ 点受到瞬时冲量 $I$ 作用，获得一个垂直于原速度方向的分速度，根据动量定理 $I = m\Delta v$

可得其大小为 $\Delta v = \frac{I}{m}$

由于冲量方向（背离地心向外，即径向）与原速度方向垂直，根据矢量合成法则（勾股定理），

作用后的合速度 $v$ 的平方为 $v^2 = v_0^2 + (\Delta v)^2 = \frac{GM}{r} + \left(\frac{I}{m}\right)^2$

则该太空碎片受到冲量作用后瞬间的动能为 $E_k = \frac{1}{2}mv^2 = \frac{1}{2}m\left[\frac{GM}{r} + \frac{I^2}{m^2}\right] = \frac{GMm}{2r} + \frac{I^2}{2m}$

故选 B。
}

\item
%% number: 8
%% paperName: 2026年河北高考第 8 题 6 分
%% typeId: 2
%% type: 多选题
%% chapter: 交流电
%% point: 交流电
%% method: 无
%% score: 6
%% degree: 550
%% duplicateId: 0
%% body:
图~\ref{2026河北08a} 为一种等离子体工作电路，理想变压器原线圈两端接电压为 $U$ 的正弦交流电源，原、副线圈的匝数分别为 $n_{1}$、$n_{2}$，$R$ 为保护电阻。稳定工作时交流电流表的读数为 $I$，若等离子体可等效为定值电阻，其两端的正弦交流电压波形如图~\ref{2026河北08b} 所示，$U_{m}$ 为电压峰值，$T$ 为周期，则下列说法正确的是\xzanswer{AD}
\twopicture[labelA=2026河北08a, labelB=2026河北08b, numA=1, numB=2, h=3.0cm]
{figs/08a.png}
{figs/08b.png}

\fourchoices[answer=AD]
{副线圈电压为 $\frac{n_{2}}{n_{1}}U$}
{原线圈中电流为 $\frac{n_{1}}{n_{2}}I$}
{电源电压周期为 $\frac{n_{2}}{n_{1}}T$}
{等离子体的等效电阻为 $\frac{\sqrt{2}U_{m}}{2I}$}

%% answer: AD
%% memo:
\memoanswer{
A．根据理想变压器原副线圈的电压比等于匝数比有 $\frac{n_{1}}{n_{2}}=\frac{U}{U'}$

可得副线圈电压为 $U' = \frac{n_{2}}{n_{1}} U$，故 A 正确；

B．根据理想变压器原副线圈的电流比等于匝数反比有 $\frac{n_{1}}{n_{2}}=\frac{I'}{I}$

可得原线圈中电流为 $I' = \frac{n_{2}}{n_{1}} I$，故 B 错误；

C．由题意可知，副线圈中电压周期为 $T$，变压器不改变交流电的周期，即电源电压周期为 $T$，故 C 错误；

D．等离子体两端电压的有效值为 $U_{\text{有}}=\frac{U_{m}}{\sqrt{2}}=\frac{\sqrt{2}U_{m}}{2}$

等离子体的等效电阻为 $R_{\text{等}}=\frac{U_{\text{有}}}{I}=\frac{\sqrt{2}U_{m}}{2I}$，故 D 正确。

故选 AD。
}

\item
%% number: 9
%% paperName: 2026年河北高考第 9 题 6 分
%% typeId: 2
%% type: 多选题
%% chapter: 直线运动
%% point: 运动的合成
%% method: 无
%% score: 6
%% degree: 450
%% duplicateId: 0
%% body:
某千斤顶的结构如图所示，四根等长杆由铰链相连。摇动手柄竖直抬升重物过程中，$A$、$C$ 两点的间距每秒均匀缩短 $2\Umm$，当 $\angle CAB = 45^\circ$ 时，下列说法正确的是\xzanswer{BD}
\onepicture[w=4.2cm, align=right]{figs/09.png}

\fourchoices[answer=BD]
{$A$ 与 $C$ 两点速度大小相等，方向相反}
{$B$ 点速度方向竖直向上，大小为 $2\Umms$}
{$C$ 点速度方向沿 $CB$ 向上，大小为 $1\Umms$}
{$A$ 点相对 $C$ 点的速度沿水平方向，大小为 $2\Umms$}

%% answer: BD
%% memo:
\memoanswer{
ABC．由题意可知，四边形 $ABCD$ 是菱形，摇动手柄竖直抬升重物过程中，$B$ 点的速度方向竖直向上，$A$、$C$ 两点的速度方向分别垂直 $AD$、$CD$ 方向，大小相等，$A$、$C$ 两点的间距每秒均匀缩短 $2\Umm$，可知 $A$、$C$ 两点速度的水平分量大小相等为 $v_{Ax} = v_{Cx} = \frac{1}{2}\times 2\Umms = 1\Umms$

当 $\angle CAB = 45^\circ$ 时，$B$ 点和 $C$ 点沿 $BC$ 方向的分速度大小相等，即 $v_B \sin 45^\circ = v_C$

又 $v_{Cx} = v_C \cos 45^\circ$

可得 $v_C = \sqrt{2}\Umms$，$v_B = 2\Umms$

$A$、$C$ 两点的速度大小相等，方向相互垂直，故 A、C 错误，B 正确；

D．$A$、$C$ 两点竖直方向的速度大小相等，方向相同，水平方向速度大小相等，方向相反，

所以 $A$ 点相对 $C$ 点的速度沿水平方向，大小为 $v_{Ax} + v_{Cx} = 2\Umms$，故 D 正确。

故选 BD。
}

\item
%% number: 10
%% paperName: 2026年河北高考第 10 题 6 分
%% typeId: 2
%% type: 多选题
%% chapter: 磁场
%% point: 安培力
%% method: 无
%% score: 6
%% degree: 350
%% duplicateId: 0
%% body:
如图所示，两平行光滑金属导轨固定在绝缘水平面上，间距为 $L$，电阻不计。左端接有输出电流大小恒为 $I$ 的恒流源，电流方向如图。导轨间分布着两个紧邻的正方形磁场区域、宽度均为 $L$，左侧磁场方向竖直向下，右侧磁场方向竖直向上，以导轨上两磁场交界处的 $O$ 点为坐标原点，沿导轨方向建立 $x$ 坐标轴，两磁场的磁感强度大小 $B$ 仅随坐标 $x$ 变化，且满足 $B = \begin{cases} 2B_{0}\frac{|x|}{L}, & 0 \le x \le \frac{L}{2} \\ 2B_{0}\left(1-\frac{|x|}{L}\right), & \frac{L}{2} < x \le L \end{cases}$ 式中 $B_{0}$ 为已知常量。$t = 0$ 时，一根质量为 $m$，导轨间电阻为 $R$ 的导体棒以初速度 $v_{0}$ 从 $x = 0$ 处向右运动；$t = t_{0}$ 时，导体棒第一次到达 $x = \frac{L}{2}$ 处，且速度为 0，运动过程中导体棒始终与导轨接触良好且垂直，不计空气阻力，忽略磁场的边界效应及电流对磁场的影响，下列说法正确的是\xzanswer{AB}
\onepicture[w=9.0cm]{\includesvg{figs/10}}

\fourchoices[answer=AB]
{导体棒做简谐运动}
{初速度 $v_{0} = \sqrt{\frac{B_{0}IL^{2}}{2m}}$}
{若初速度变为 $0.5v_{0}$，则导体棒的速度第一次变为 $0$ 所需时间为 $\frac{\sqrt{2}}{2}t_{0}$}
{若初速度变为 $1.5v_{0}$，则导体棒将以两磁场交界线为中心做往复运动}

%% answer: AB
%% memo:
\memoanswer{
A．作出 $B-x$ 图像如图所示

\onepicture[w=5.0cm]{figs/10_an.png}

根据题意知，在 $0 \leq x \leq \frac{L}{2}$ 区域，导体棒所受安培力

大小 $F_{\text{安}}=BIL=2B_{0}Ix$

根据左手定则可知安培力方向始终指向 $x=0$ 处，令 $2B_{0}I = k$，故导体棒所受外力满足 $F = -kx$。

同理，导体棒在 $-\frac{L}{2} \leq x \leq 0$ 区域也满足 $F = -kx$

则导体棒在 $-\frac{L}{2} \leq x \leq \frac{L}{2}$ 区域做简谐运动，故 A 正确；

B．导体棒 $x = 0$ 到 $x = \frac{L}{2}$ 过程，根据动能定理 $-\overline{F_{\text{安}}} \cdot \frac{L}{2} = 0 - \frac{1}{2}mv_{0}^{2}$

其中 $\overline{F_{\text{安}}} = \frac{0 + 2B_{0}IL}{2} = \frac{B_{0}IL}{2}$

联立可得 $v_{0} = \sqrt{\frac{B_{0}IL^{2}}{2m}}$，故 B 正确；

C．对于简谐运动，从平衡位置运动到最大位移处所用时间为 $t = \frac{T}{4} = \frac{\pi}{2\omega}$

根据简谐运动圆频率 $\omega = \sqrt{\frac{|k|}{m}}$

结合上述分析可得 $\omega = \sqrt{\frac{2B_{0}I}{m}}$

$\omega$ 与速度无关，初速度变为 $0.5v_{0}$ 后，根据动能定理可知导体棒向右运动的最大距离小于 $x = \frac{L}{2}$，导体棒仍做简谐运动，根据上述分析可知速度变为 $0.5v_{0}$ 对导体棒的速度第一次变为 $0$ 所需的时间没有影响，仍为 $t_{0}$，故 C 错误；

D．若初速度变为 $1.5v_{0}$，由动能定理 $-\overline{F_{\text{安}}} \cdot \frac{L}{2} = \frac{1}{2}mv^{2} - \frac{1}{2}m(1.5v_{0})^{2}$

可得导体棒运动到 $x = \frac{L}{2}$ 处速度为 $v = \frac{\sqrt{5}}{2}v_{0} > v_{0}$

根据对称性知导体棒从 $x = \frac{L}{2}$ 处运动到 $x = L$ 处克服安培力做功为 $W = \overline{F_{\text{安}}}\cdot\frac{L}{2} = \frac{2B_{0}I\frac{L}{2} + 0}{2}\cdot\frac{L}{2} = \frac{1}{2}mv_{0}^{2} < \frac{1}{2}mv^{2}$

即导体棒运动到 $x = L$ 处速度不为零，导体棒将继续向右运动离开磁场，故导体棒不会以两磁场交界线为中心做往复运动，故 D 错误。

故选 AB。
}

\item
%% number: 11
%% paperName: 2026年河北高考第 11 题 4 分
%% typeId: 4
%% type: 实验题
%% chapter: 电场
%% point: 电容器充放电实验
%% method: 无
%% score: 4
%% degree: 650
%% duplicateId: 0
%% body:
\begin{enumerate}
\item
某同学用弹簧测力计探究作用力和反作用力的关系。将甲、乙两测力计水平连接在一起，测力左侧固定，用手向右缓慢拉测力计乙。拉至某位置时两测力计示数如图所示。
\onepicture[w=9.5cm, align=right]{figs/11a.png}
\begin{steps}
\item
测力计乙的读数是 \tkanswer{1.98} $\UN$。
\item
实验中发现两测力计示数总是明显不同、原因可能是 \tkanswer{C}。（单选）
\threechoices[answer=C]
{两测力计的量程不同}
{两测力计弹簧的材质不同}
{使用测力计前未正确调零}
\end{steps}
\item
如图为探究电容器充、放电实验的电路图，实验选用内阻为几千欧姆的电压表和内阻很小且零刻度的电流表。
\onepicture[w=5.0cm, align=right]{figs/11b.png}
\begin{steps}
\item
实验时将开关 $S$ 接 1 端，可观察到电流表指针 \tkanswer{B}（填字母，“A．迅速偏转到某位置后稳定不动”或“B．迅速偏转到某位置后往回偏转”），同时电压表指针偏转到某位置后稳定不动。
\item
现将 $S$ 与 1 端断开，但尚未与 2 端连接，此时发现电压表的指针逐渐向零刻度偏转。经检查器材完好，线路连接正确，出现此现象的原因可能是 \tkanswer{开关断开时电容器相当于和电压表形成了闭合回路，电压表内阻较大，放电电流较小，电容器放电过程中两端电压逐渐减小直至放电结束}。
\end{steps}
\end{enumerate}

%% answer: （1）1.98（1.96～2.00 均可），C；（2）B，电容器通过电压表放电
\jdanswer{
\begin{enumerate}
\item
$1.98$（$1.96 \sim 2.00$ 均可），C
\item
B，开关断开时电容器相当于和电压表形成了闭合回路，电容器通过电压表缓慢放电
\end{enumerate}
}

%% memo:
\memoanswer{
（1）①由图可知弹簧测力计的分度值为 $0.1\UN$，需要往后估读一位，因此读数为 $1.98\UN$。

②探究作用力和反作用力的关系实验中，根据牛顿第三定律可知作用力和反作用力总是大小相等，量程和弹簧测力计弹簧的材质都不影响其示数的准确性，两个测力计示数总是明显不同原因可能是测力计未正确调零。故选 C。

（2）①电容器充电过程中，开关闭合时充电电流瞬间达到较大的值，然后又逐渐减小，充电完成时，电流为 $0$，因此观察到的现象是电流表指针迅速偏转到某位置后往回偏转。故选 B。

②观察题图可以发现，开关断开时电容器相当于和电压表形成了闭合回路，电压表内阻为几千欧姆，比较大，因此放电电流比较小，电容器放电的过程中其两端电压逐渐减小，直至最终放电结束，电容器带电荷量变为 $0$，电压表示数减为 $0$。
}

\item
%% number: 12
%% paperName: 2026年河北高考第 12 题 8 分
%% typeId: 4
%% type: 实验题
%% chapter: 电路
%% point: 传感器
%% method: 无
%% score: 8
%% degree: 550
%% duplicateId: 0
%% body:
问题探究

某学生兴趣小组通过查阅资料了解到，利用电磁阀控制压缩空气驱动活塞运动可实现公交车门开合，遂以“电控气动技术的原理探索与应用”为主题开展探究式学习。
\begin{enumerate}
\item
分析电控气动门工作原理。

图~\ref{2026河北12a} 是电控气动系统开门过程的简化图。电磁阀中的固定铁芯与线圈组成电磁铁，当开关 $S_{1}$ 闭合时，衔铁在电磁铁的吸引下带动三个阀芯一起向左运动，复位弹簧被压缩，储气罐内压缩空气经气孔 B 进入汽缸右侧气室。推动汽缸内活塞向 \tkanswer{左}（填“左”或“右”）运动，车门打开。同时，汽缸左侧气室中的气体从气孔 \tkanswer{R}（填“S”或“R”）排出。当 $S_{1}$ 断开时，复位弹簧推动衔铁和阀芯向右运动，压缩空气推动活塞，车门关闭。
\onepicture[num=1, label=2026河北12a, w=6.0cm, align=right]{figs/12a.png}

\item
探究复位弹簧的受力情况。

同学们测量了复位弹簧在弹性限度内受到的作用力 $F$ 与形变量 $x$ 的数据，并标在了坐标纸上，如图~\ref{2026河北12b} 所示，请在答题卡上绘制出该弹簧的 $F-x$ 图像；根据绘制的 $F-x$ 图像，得到该弹簧的劲度系数 $k =$ \tkanswer{0.22} $\UNmm$（保留两位有效数字）。

同学们据此估算出电磁阀中电磁铁对阀芯的作用力仅有几牛顿，而资料显示，公交车实现车门开合所需的作用力高达上千牛顿。系统利用电控气动技术实现了“四两拨千斤”的效果。
\onepicture[num=2, label=2026河北12b, w=5.0cm, align=right]{figs/12b.png}

\item
电控气动技术的应用设计。

受此启发，同学们设计了一款温室自动开窗器，电路如图~\ref{2026河北12c} 所示。控制电路板的程序能够实现温度高于 $35\UdC$ 时自动闭合开关 $S_{1}$，温度低于 $25\UdC$ 时自动断开 $S_{1}$。电磁阀线圈工作电压为 $24\UV$，而控制电路板的工作电压（$5\UV$）和电流较小，需配置电磁继电器作为控制开关。请在答题卡相应图上连线、完成电路、实现如下效果：$S_{2}$ 闭合时，电磁阀中的电磁铁吸引衔铁左移，窗户打开；$S_{2}$ 断开时，电磁阀中的衔铁右移，窗户关闭。
\onepicture[num=3, label=2026河北12c, w=8.0cm, align=right]{figs/12c.png}
\end{enumerate}

%% answer: （1）左，R；（2）0.22；（3）连线见解析
\jdanswer{
\begin{enumerate}
\item 左，R
\item $0.22\UNmm$
\item 连线如解析图所示
\end{enumerate}
}

%% memo:
\memoanswer{
（1）结合题意可知，右侧气室充入气体，若体积不变，气体压强将会增大，推动活塞向左运动；由题图可知，左侧气室出气孔与 R 相连，则气体从 R 排出。

（2）作图时要注意到图线要通过尽可能多的点画出直线，不在直线上的点应该均匀分布在直线的两侧，$F-x$ 图像如图所示

\onepicture[w=5.0cm]{figs/12_an1.png}

根据 $F = kx$ 可知 $F - x$ 图像的斜率即为劲度系数，得 $k = \frac{2.5}{11.5}\UNmm \approx 0.22\UNmm$。

（3）$S_{2}$ 闭合时，电磁继电器的线圈内有了电流，会使与 3 和 2 相连的两个触点接触，并使 K、J 连通，右侧电磁阀电路形成闭合回路，衔铁左移，窗户打开；$S_{2}$ 断开时，电磁继电器的线圈内失去电流，与 3 和 2 相连的两个触点分开，导致 K、J 部分是断路状态，电磁阀部分的衔铁在复位弹簧作用下右移、窗户关闭，连线情况如图所示

\onepicture[w=8.0cm]{figs/12_an2.png}
}

\item
%% number: 13
%% paperName: 2026年河北高考第 13 题 8 分
%% typeId: 6
%% type: 计算题
%% chapter: 机械波
%% point: 波的多解性
%% method: 无
%% score: 8
%% degree: 550
%% duplicateId: 0
%% body:
下图为一列沿 $x$ 轴正方向传播的简谐横波在 $t = 0$ 时的波形（中间部分未画出），此时平衡位置在 $x = x_{P}$ 的质点处于波峰。
\begin{enumerate}
\item
写出该波的波长，以及 $x_{P}$ 在 $10 \sim 14\Um$ 范围内的所有可能值。
\item
若平衡位置位于原点处的质点在 $0 \sim 1\Us$ 内通过的路程为 $5\Ucm$，求该波的周期和波速。
\end{enumerate}
\onepicture[w=10.0cm, align=right]{figs/13.png}

%% answer:
\jdanswer{
\begin{enumerate}
\item
$2\Um$，$11\Um$ 和 $13\Um$

\item
$0.8\Us$，$2.5\Ums$
\end{enumerate}
}

%% memo:
\memoanswer{
（1）相邻波峰（或波谷）间距等于波长，根据 $t = 0$ 时的波形图可知该波的波长为 $\lambda = 2\Um$。

由图可知，平衡位置处于波峰的质点 $x = x_{P}$ 满足 $x_{P} = \left(n + \frac{1}{2}\right)\lambda(n = 2, 3, 4, \cdots)$

又 $10\Um < x_{P} < 14\Um$

$x_{P}$ 在 $10 \sim 14\Um$ 范围内的可能值有：$n = 5$ 时，$x_{P} = 11\Um$；

$n = 6$ 时，$x_{P} = 13\Um$。

$x_{P}$ 在 $10 \sim 14\Um$ 范围内的所有可能值有 $11\Um$ 和 $13\Um$。

（2）由题意可知，振幅 $A = 1\Ucm$，$s = 5A = 5\Ucm$

则有 $1\Us = T + \frac{1}{4}T$

由波速公式有 $v=\frac{\lambda}{T}$

解得 $T = 0.8\Us$，$v = 2.5\Ums$。
}

\item
%% number: 14
%% paperName: 2026年河北高考第 14 题 14 分
%% typeId: 6
%% type: 计算题
%% chapter: 运动和力
%% point: 动量和能量
%% method: 无
%% score: 14
%% degree: 450
%% duplicateId: 0
%% body:
如图所示，质量为 $4\Ukg$ 的木板上放有一个质量为 $1\Ukg$ 的机器人，木板始终受到水平向右、大小为 $5\UN$ 的恒力作用。初始时木板与机器人一起以 $1\Ums$ 的速度沿水平地面向右匀速运动。机器人正上方有一个沿竖直方向可以伸缩、水平向右速度恒为 $1\Ums$ 的机械夹爪。某时刻夹爪将机器人向上提起，$2\Us$ 后放回木板，同时夹爪缩回，机器人在摩擦力的作用下最终与木板相对静止。$g$ 取 $10\Umsq$，机器人可视为质点，机器人被提起和放下瞬间竖直方向速度均为零。求
\begin{enumerate}
\item
机器人被提起的 $2\Us$ 内，木板位移的大小。
\item
从机器人被放回木板到与木板相对静止的过程中，摩擦力对机器人所做的功。
\end{enumerate}
\onepicture[w=9.0cm, align=right]{figs/14.png}

%% answer:
\jdanswer{
\begin{enumerate}
\item
$2.5\Um$

\item
$0.48\UJ$
\end{enumerate}
}

%% memo:
\memoanswer{
（1）根据题意，设木板与地面间的摩擦因数为 $\mu$，则有 $F = \mu \left( m_{\text{木}} + m_{\text{机}} \right) g$

解得 $\mu = 0.1$

机器人被提起时，对木板有 $F - \mu m_{\text{木}}g = m_{\text{木}}a_{1}$

解得 $a_{1}=0.25\Umsq$

机器人被提起的 $2\Us$ 内，木板位移的大小 $x = v_{0}t + \frac{1}{2}a_{1}t^{2} = \left(1 \times 2 + \frac{1}{2} \times 0.25 \times 2^{2}\right)\Um = 2.5\Um$

（2）机器人被放回木板时，木板的速度为 $v_{2}=v_{0}+a_{1}t=(1+0.25\times2)\Ums=1.5\Ums$

机器人被放回木板后，恒力与地面对木板的摩擦力平衡，机器人和木板组成的系统所受合力为零，则由动量守恒定律有 $m_{\text{机}}v_{0}+m_{\text{木}}v_{2}=\left(m_{\text{机}}+m_{\text{木}}\right)v$

解得 $v = 1.4\Ums$

对机器人，由动能定理可得，摩擦力对机器人所做的功 $W = \frac{1}{2} m_{\text{机}} v^{2} - \frac{1}{2} m_{\text{机}} v_{0}^{2} = 0.48\UJ$。
}

\item
%% number: 15
%% paperName: 2026年河北高考第 15 题 16 分
%% typeId: 6
%% type: 计算题
%% chapter: 磁场
%% point: 带电粒子在组合场中的运动
%% method: 无
%% score: 16
%% degree: 350
%% duplicateId: 0
%% body:
图~\ref{2026河北15a} 为某质谱仪工作原理示意图。电离室中的气体分子被激光照射后发生电离，其中带正电的粒子静止经平行于纸面的加速电场加速后，垂直于 $CD$ 边进入梯形匀强磁场区域（磁场方向垂直纸面向外），并且 $CG$ 边中点 $O$ 平行于 $CD$ 边射出，经无场区从边界 $PQ$ 进入平行于纸面的匀强偏转电场，最终打到接收器与纸面交线 $MN$ 上并被吸收，接收器可视为接地良好的金属板。$MN$ 延长线经过 $O$ 点，与 $CD$ 所在直线夹角为 $\alpha$，$\alpha$ 可通过 $MN$ 绕 $O$ 点在纸面内的转动进行调整，调整前后 $PQ$ 与 $MN$ 始终平行且间距为 $0.3\Um$ 不变。$CG$ 边长 $0.2\Um$，与底边 $HG$ 夹角为 $60^\circ$，磁感应强度大小为 $0.4\UT$，偏转电场的电场强度大小为 $1.8\times10^{4}\UVm$，方向始终垂直于 $PQ$。整个过程只考虑一种粒子，加速电场的电压恒定，$MN$ 与 $PQ$ 足够长，装置处于真空环境，忽略粒子间的相互作用，不计重力。
\begin{enumerate}
\item
已知激光波长为 $442\Unm$，求该激光一个光子的能量 $\varepsilon$。（普朗克常量 $h = 6.63\times10^{-34}\UJcdots$，真空中光速 $c = 3\times10^{8}\Ums$）
\item
调整 $\alpha$，得到从粒子离开加速电场到到达接收器所经历的时间 $t$ 与 $\alpha$ 的关系如图~\ref{2026河北15b} 所示，求加速电场的电压 $U$ 以及粒子比荷 $k$。
\item
质谱仪稳定工作时，测得接收器每秒接收的粒子数为 $n$，若再测一个除电荷量和质量以外的物理量，便可在（2）问的基础上得到粒子的质量。请写出该物理量，并推导粒子质量表达式（所有物理量均用字母表示）。
\end{enumerate}
\twopicture[labelA=2026河北15a, labelB=2026河北15b, numA=1, numB=2, hA=4.5cm, hB=3.8cm, align=right]
{figs/15a.png}
{figs/15b.png}

%% answer:
\jdanswer{
\begin{enumerate}
\item
$\varepsilon = 4.5\times10^{-19}\UJ$

\item
$k = 6\times10^{6}\UCkg$，$U = 3600\UV$

\item
测接收器受到粒子的平均作用力大小 $F$，质量表达式为 $m = \frac{F}{n\sqrt{2kU}\sin\alpha}$；测流过地线上的平均电流 $I$，质量表达式为 $m = \frac{I}{nk}$
\end{enumerate}
}

%% memo:
\memoanswer{
（1）根据光子能量公式 $\varepsilon = h\nu = \frac{hc}{\lambda}$

代入数据解得 $\varepsilon = 4.5 \times 10^{-19}\UJ$。

（2）根据题意画出粒子运动轨迹图如下

\onepicture[w=6.0cm]{figs/15_an.png}

根据几何关系可知粒子在磁场半径 $R = 0.05\sqrt{3}\Um$

且粒子在磁场中洛伦兹力提供向心力，有 $qvB = m\frac{v^{2}}{R}$

其中比荷 $k=\frac{q}{m}$。可得 $v=kBR=0.02\sqrt{3}k$

由题知粒子在无场区 $O$ 到 $PQ$ 的距离为 $\frac{d}{\sin\alpha}$，在无场区运动时间 $t_2 = \frac{2d}{v\sin\alpha}$

粒子进入电场时的法向初速度与电场方向相反，大小为 $v_y = v\sin\alpha$

则在法线方向有 $t_3 = \frac{2v_y}{a}$，$a = \frac{qE}{m} = kE$

解得 $t_{3}=\frac{2v\sin\alpha}{kE}$

总时间 $t=t_{0}+t_{1}+t_{2}+t_{3}$

其中 $t_{0}$ 为加速电场末端到 $CD$ 无场区的时间与 $\alpha$ 无关，$t_{1}=\frac{\pi m}{2qB}=\frac{\pi}{2kB}$ 与 $\alpha$ 无关，根据图~\ref{2026河北15b} 可知

$\alpha=60^{\circ}$ 时 $t$ 最小，即 $t'=\frac{2d}{v\sin\alpha}+\frac{2v\sin\alpha}{kE}$

$t'$ 在 $\alpha = 60^\circ$ 有最小值，根据均值不等式可求出 $v^2 = \frac{2}{3} k Ed$

联立 $v = kBR$

代入数据得 $k = 6 \times 10^{6}\UCkg$

在加速电场中有 $\frac{1}{2}mv^{2}=qU$

得 $U=\frac{v^{2}}{2k}=3600\UV$

（3）方法一：测量接收器受到粒子的平均作用力大小 $F$，在 $\Delta t$ 时间内，有 $N = n\Delta t$

根据动量定理和牛顿第三定律可知接收器受到的平均作用力大小 $F$ 等于粒子动量的变化率大小，即 $F = \frac{N\Delta p}{\Delta t} = n\Delta p$

对单个粒子，设粒子打在接收器上时垂直接收器速度为 $v_y$，则有 $\Delta p = mv_y$

可得 $m = \frac{F}{nv_y}$

根据（2）问可知 $v_{y}=v\sin\alpha$

联立可得粒子质量 $m = \frac{F}{n\sqrt{2kU}\sin\alpha}$

方法二：已知每秒接收粒子数 $n$，每个粒子电荷量为 $q$，则电流 $I = nq$

故粒子质量 $m = \frac{q}{k} = \frac{I}{nk}$

将 $k$ 代入即可。
}

\end{enumerate}

\end{document}
